%% LyX 2.5.1 created this file.  For more info, see https://www.lyx.org/.
%% Do not edit unless you really know what you are doing.
\documentclass[english]{article}
\usepackage[T1]{fontenc}
\usepackage[cp862]{inputenc}
\usepackage{xcolor}
\usepackage{amsmath}
\usepackage{amsthm}
\usepackage{amssymb}
\usepackage{geometry}
\geometry{verbose,tmargin=0.07\textwidth,bmargin=0.1\textwidth,lmargin=0.1\textwidth,rmargin=0.1\textwidth}
\PassOptionsToPackage{normalem}{ulem}
\usepackage{ulem}

\makeatletter

%%%%%%%%%%%%%%%%%%%%%%%%%%%%%% LyX specific LaTeX commands.
\providecolor{lyxadded}{rgb}{0,0,1}
\providecolor{lyxdeleted}{rgb}{1,0,0}
%% Change tracking with ulem and xcolor: base macros
\DeclareRobustCommand{\mklyxadded}[1]{\textcolor{lyxadded}\bgroup#1\egroup}
\DeclareRobustCommand{\mklyxdeleted}[1]{\textcolor{lyxdeleted}\bgroup\mklyxsout{#1}\egroup}
\DeclareRobustCommand{\mklyxsout}[1]{\ifx\\#1\else\sout{#1}\fi}
%% Change tracking with ulem and xcolor: ct markup
\DeclareRobustCommand{\lyxadded}[4][]{\mklyxadded{#4}}
\DeclareRobustCommand{\lyxdeleted}[4][]{\mklyxdeleted{#4}}

%%%%%%%%%%%%%%%%%%%%%%%%%%%%%% Textclass specific LaTeX commands.
\numberwithin{equation}{section}
\numberwithin{figure}{section}

%%%%%%%%%%%%%%%%%%%%%%%%%%%%%% User specified LaTeX commands.
\usepackage{algorithm2e}
\usepackage{binarytree}
\usepackage{tikz}
\usepackage{pgfplots}
\pgfplotsset{compat=newest}
\usepackage{amsmath}
\DeclareMathOperator{\sgn}{sgn}
\newcommand{\lcm}[1]{\operatorname{lcm}\left(#1\right)}

\makeatother

\usepackage{babel}
\begin{document}
\title{NLP-Assignment 1}
\author{Yonatan Karni | Raz Ziv}
\maketitle

\section*{3.}

\subsection*{a.}

Firstly we notice that for fixed $o,c$ the number of times $p_{\theta}\left(o|c\right)$
shows in the objective is $\#\left(c,o\right)$,

This follows from the fact that in each window where both $o,c$ is
in we count $p_{\theta}\left(o|c\right)$ one more time.

\[
\mathcal{L}\left(p_{\theta}\left(c|o\right)\right)=\#\left(c,o\right)\log\left(p_{\theta}\left(o|c\right)\right)
\]

Then we can rewrite the objective as:

\[
\mathcal{L}\left(\theta\right)=\sum^{T}_{t}\sum_{-m\le j\le m}\log\left(p_{\theta}\left(o|c\right)\right)=\sum_{c\in Voc}\sum_{o\in Voc}\#\left(c,o\right)\log\left(p_{\theta}\left(o|c\right)\right)=\sum_{c\in Voc}\sum_{o\in Voc}\mathcal{L}\left(p_{\theta}\left(c|o\right)\right)
\]

For a fixed $o,c$ we must make sure $\sum_{o'\in Voc}p_{\theta}\left(o'|c\right)=1$
namely $\theta$ is a probability distribution.

Therefore we will use Lagrange multiplier with $\lambda$ parameter
to enforce this constraint.

Moreover, we can see that maximizing the original objective $\mathcal{L}$
is equivalent to maximize each term in the sum over the words $o,c$. 

\begin{align*}
\theta^{*} & =\arg\max_{\theta}\mathcal{L}\left(\theta\right)\\
\iff & \forall\left(c,o\right)\ p_{\theta^{*}}\left(o|c\right)=\arg\max_{p_{\theta}\left(c|o\right)}\mathcal{L}\left(p_{\theta}\left(c|o\right),\lambda\right)\\
s.t\ \mathcal{L}\left(p_{\theta}\left(c|o\right),\lambda\right)= & \left(\#\left(c,o\right)\right)\log\left(p_{\theta}\left(o|c\right)\right)+\lambda\left(1-\sum_{o'\in Voc}p_{\theta}\left(o'|c\right)\right)
\end{align*}

Differentiating by $p_{\theta}\left(o|c\right)$ and comparing to
$0$ we get:

\begin{align*}
\frac{\partial\mathcal{L}\left(p_{\theta}\left(c|o\right),\lambda\right)}{\partial p_{\theta}\left(c|o\right)} & =\frac{\#\left(c,o\right)}{p_{\theta^{*}}\left(o|c\right)}-\lambda=0\\
p_{\theta^{*}}\left(o|c\right) & =\frac{\#\left(c,o\right)}{\lambda}
\end{align*}

Now substitute back into the original Lagrange multiplier definition:

\begin{align*}
\sum_{o'\in Voc}p_{\theta^{*}}\left(o'|c\right) & =\sum_{o'\in Voc}\frac{\#\left(c,o'\right)}{\lambda}=1\\
\lambda & =\sum_{o'\in Voc}\#\left(c,o'\right)
\end{align*}

Finally we obtain:

\[
p_{\theta^{*}}\left(o|c\right)=\frac{\#\left(c,o\right)}{\sum_{o'\in Voc}\#\left(c,o'\right)}
\]


\subsection*{b.}

We will take the following corpus $\left\{ abc,cbc,aac\right\} $
over the alphabet $\left\{ a,b,c\right\} $

Using a window of size 3 for skip-gram model.

Let us assume $p$ is optimal while being expressed with the given
softmax expression.

Lets calculate the optimal probabilities using the optimal form we
found in last section:

\begin{align*}
p\left(a|c\right) & =\frac{3}{8},\ p\left(a|a\right)=\frac{2}{6}\\
p\left(b|c\right) & =\frac{3}{8},\ p\left(b|a\right)=\frac{1}{6}\\
p\left(c|c\right) & =\frac{2}{8},\ p\left(c|a\right)=\frac{3}{6}
\end{align*}

Now we can obtain the following on the embedding of $v_{c},v_{a}$
:

\begin{align*}
\log\left(\frac{p\left(a|c\right)}{p\left(c|c\right)}\right) & =\log\left(\frac{e^{u_{a}v_{c}}}{e^{u_{c}v_{c}}}\right)=v_{c}\left(u_{a}-u_{c}\right)=\log\left(\frac{3}{2}\right)\\
\log\left(\frac{p\left(a|a\right)}{p\left(c|a\right)}\right) & =\log\left(\frac{e^{u_{a}v_{a}}}{e^{u_{c}v_{a}}}\right)=v_{a}\left(u_{a}-u_{c}\right)=\log\left(\frac{2}{3}\right)\\
\iff & \frac{v_{c}}{v_{a}}=\frac{\log\left(\frac{3}{2}\right)}{\log\left(\frac{2}{3}\right)}
\end{align*}

On the other hand:

\begin{align*}
\log\left(\frac{p\left(b|c\right)}{p\left(c|c\right)}\right) & =\log\left(\frac{e^{u_{b}v_{c}}}{e^{u_{c}v_{c}}}\right)=v_{c}\left(u_{b}-u_{c}\right)=\log\left(\frac{3}{2}\right)\\
\log\left(\frac{p\left(b|a\right)}{p\left(c|a\right)}\right) & =\log\left(\frac{e^{u_{b}v_{a}}}{e^{u_{c}v_{a}}}\right)=v_{a}\left(u_{b}-u_{c}\right)=\log\left(\frac{1}{3}\right)\\
\iff & \frac{v_{c}}{v_{a}}=\frac{\log\left(\frac{3}{2}\right)}{\log\left(\frac{1}{3}\right)}
\end{align*}

From both we get the following contradiction:

\[
\frac{\log\left(\frac{3}{2}\right)}{\log\left(\frac{2}{3}\right)}=\frac{\log\left(\frac{3}{2}\right)}{\log\left(\frac{1}{3}\right)}\iff\log\left(3\right)=\log\left(2\right)
\]

\pagebreak

\section*{5.}

\subsection*{5.1.a.}

I will discuss two intrinsic evaluation metrics:

\textbf{Coherence-}

measure how much a certain topic is coherent.

Namely, how much top words that belong to the topic are likely to
be seen near each other in the corpus.

For example, given the following words inside a topic:

``surf sea wave shore'' will give higher coherence then ``surf
champion banana car''

simply because the first words are more likely to appear in the same
context.\\

The mathematical expression for the coherence:

First we will choose some top matching words $W$ for a certain topic
$T$.

Denote $P\left(w\right)$ the probability to get the word $w$ from
some context, and $P\left(u,w\right)$ is the joint distribution ,
ie to see both words in the same context.

\begin{align*}
\text{Coherence}\left(T\right) & =\frac{1}{\left|W\right|}\sum_{\left(w,u\right)\in W}\log\left(\frac{P\left(w,u\right)}{P\left(w\right)P\left(u\right)}\right)\\
\text{Coherence} & =\frac{1}{K}\sum^{K}_{i=1}\text{Coherence}\left(T_{i}\right)
\end{align*}

\textbf{Diversity-}

The diversity metric aims to measure how much the topics differ.

We measure how likely that a word will appear in both topics in the
following method:

Choose $W_{k}$ set of most likely words for each topic.

Assume $K$ topics and $N$ top words per topic.

\[
\textnormal{Diversity}=\frac{\left|\bigcup^{K}_{k=1}W_{k}\right|}{K\cdot N}
\]


\subsection*{5.1.b.}

I saw several topic models use wikipedia as a benchmark dataset for
evaluation. 

20 Newsgroup, a large dataset contains news posts on 20 topics is
used in different research papers.

Trump's Twitter messages archive was used in BERTopic evaluation git
repository.

\subsection*{5.3.}

I used the two coherence parameters $c\_v$ and $c\_uci$.

Both parameters use a sliding window to check if some words appear
together, $c\_v$ is more sophisticated with additional statistical
methods and different normalization.

I got the following results:

\begin{align*}
c\_v & =0.41\\
c\_uci & =-5.97
\end{align*}

The first parameter range $0\le c\_v\le1$ where $0$ means no coherence
at all and $1$ means very strong coherence. We got $0.41$ so we
can see a coherence was detected but not that high, probably due to
real world messy text data or not enough data.

The second parameter $c\_uci\in\mathbb{R}$ however most of times
it falls approximately in $[-15,5]$.

It shows a medium coherence measure which probably results from its
strict window of length $10$, some words appear for example in the
next sentence.

\subsection*{5.4}

One problem is semantics, namely the literal meaning of words. The
two evaluation parameters I have shown earlier can fail on this regard. 

For example a topic may contain top words which are written close
to each other many times and therefore produce high coherence, but
the words are very different semantically. \\

Another example is a model with very low diversity which may seem
bad. But a human may recognize a lot of topics being merged together
in the corpus, for example science of music. 
\end{document}
